\section{Remembering the Dead}

\begin{quotex}
This translation of Charles Maurras' \textit{All Soul's Day}\footnote{\url{http://www.gornahoor.net/library/AllSoulsDay.pdf}} originally appeared in Taki's Magazine\footnote{\url{http://takimag.com/article/all\_souls\_day/}}.
\end{quotex}

On \textbf{All Soul's Day}, we honor the memory of the dead, especially those closest to us in our hearts and mind. \textbf{Charles Maurras} refers to this as ``the universal cult of the dead, of all the dead, of those who had existed, provided that they had belonged to the human race''; this is a subjective memory. Although it is praiseworthy to remember our ancestors, Maurras points out this drawback:

\begin{quotex}
in teaching us to venerate all of defunct humanity, it trains us logically to venerate, en bloc, all of living humanity, that is to say, to make us accept and even venerate the worse faults that it commits even though we recognize them as much in ourselves as in our neighbours.
\end{quotex}


Alongside this, there is an objective memory in the sense that it is part of our experience and, therefore, quite real in the Positivist sense. Maurras describes this cult as:

\begin{quotex}
more reserved, prouder, and, in my sense, more beneficial. That renders to the elite among the dead, those whom the positivists call, a little verbosely, ``the great types of humanity'', and the Catholics, more tersely, the ``saints''.
\end{quotex}


This form of remembrance has a certain benefit:

\begin{quotex}
by obliging us to hold the dead as our models, it forces us to select from among these scattered people, hence, indirectly, to make a critique of our own characters: by applying our minds to consider those great dead men, it opens us up to the way of personal exaltation and perfection.
\end{quotex}


Maurras summarizes:

\begin{quotex}
The consequence is that human solidarity, to use that term, must belong much less to the crowd of our predecessors, than to the persons of the past who have realized, in a great way, the fine natural traits of man. Those who pass up the opportunity to serve their great memory, pass up an undoubted opportunity to help themselves, to correct themselves, and to improve themselves.
\end{quotex}


Comte's follower, \textbf{Antoine Baumann}, expands on this idea:

\begin{quotex}
\emph{Humanity} does not mean in any way the whole of men scattered among the living on this planet, nor the simple total of the living and the dead. It is only the group of men who have cooperated in the great human work, those who persist in us, whom we continue, those of whom we are the true debtors, the others being at times only ``parasites'' or ``producers of manure''. This large human elite is not a futile image. It forms what there is of the most real in us. We feel it as soon as we reach down to the secret of our nature. Subjects of mathematical and astronomical facts, subjects of physical facts, chemical facts and the facts of life, we are moreover subjects of special facts in the human family. We depend on our contemporaries. We depend even more on our predecessors. What thinks in us, before us, is the human language, which is not our personal work, but the work of humanity, it is also human reason which preceded us, which surrounds us and anticipates us; it is human civilisation in which a personal contribution, as powerful that it may be, is never but a molecule of a tiny energy in the drop of water added by our contemporaries in the current of this vast river. Actions, thoughts, or feelings, are products of the human soul: our personal soul is there almost for nothing. The true positivist repeats rather like Saint Paul: in ea vivimus, movemur et sumus ``in him we live, move and have our being'' , and if he put his heart in harmony with his science and his faith, he can only add, in an act of adoration, the slightly modified words from the Psalmist: Non nobis, Domine, non nobis, sed nomini tuo da gloriam! ``Not To Us O Lord Not To Us But To Your Name Give Glory'' 
\end{quotex}


\textbf{Auguste Comte} explains the objective reality of the \textbf{Great Being}:


\begin{quotex}
The most important applications constantly derive from the theories formed in a simple scientific purpose and which often were cultivated during several centuries without producing any practical result. One can cite a very remarkable example of it in the beautiful speculations of Greek geometers on conic sections which, after a long series of generations, served, in determining the renovation of astronomy, to finally lead to the art of navigation to the degree that it attained in these days and to which it would never have appeared without the purely theoretical works of Archimedes and Apollonius; so that Condorcet was able to say with reason in this matter: ``The sailor whom an exact observation of longitude preserves from the shipwreck owes his life to a theory conceived two thousand years before, by men of genius who had in view simple geometric speculations.''
\end{quotex}


This is the time to bring to mind those ``persons of the past'' who have, over the course of time, built up the edifice of civilisation, not because they need it, but rather because we do. These persons are not always those who receive the most public adulation, which is often a mask for contemporary ideology. Instead, seek out the neglected and unknown, those whose ideas still live in the Great Being, though they may be suppressed in our time. The religious have their saints and mystics; the knights had the cult of the Nine Worthies\footnote{\url{http://www.gornahoor.net/?p=331}}; Hermetists recall their own predecessors. There are forgotten men of genius who still affect our lives; do not neglect them.


\flrightit{Posted on 2010-11-01 by Cologero}

